\documentclass[11pt]{article}
\usepackage[T1]{fontenc}
\usepackage{amsmath,amssymb,amsthm}
\usepackage[margin=1in]{geometry}
\usepackage[hidelinks]{hyperref}
\hypersetup{
  pdftitle={The Four-Color Game on M2(C3): An Exact Strategy Certificate},
  pdfauthor={Xiaojian Zeng},
  pdfsubject={Game chromatic number of a generalized Mycielski graph},
  pdfkeywords={graph coloring game, game chromatic number, generalized Mycielski graph, exact certificate}
}

\newtheorem{theorem}{Theorem}
\newtheorem{lemma}[theorem]{Lemma}

\newcommand{\gamechi}{\chi_{\mathrm g}}

\title{The Four-Color Game on $M_2(C_3)$:\\An Exact Strategy Certificate}
\author{Xiaojian Zeng\\
PBC School of Finance, Tsinghua University, 100190\\
\texttt{xiaojianzeng@foxmail.com}}
\date{}

\begin{document}
\maketitle

\begin{abstract}
We determine the game chromatic number of the ten-vertex generalized
Mycielski graph $M_2(C_3)$.  An elementary forced-color argument gives the
lower bound four.  For the upper bound, we give a complete machine-readable
Alice strategy for the four-color game.  After quotienting by simultaneous
permutations of the three indices and by permutations of the four colors, the
strategy is a directed acyclic graph with $1020$ states and $2467$ edges.  A
separately implemented verifier reconstructs the graph and the full symmetry
action, checks the strategy move at every Alice state, and checks every raw
legal reply at every Bob state.  Backward induction on the number of uncolored
vertices therefore proves
\[
  \gamechi\bigl(M_2(C_3)\bigr)=4.
\]
\end{abstract}

\section{Context and statement}

In the vertex-coloring game, Alice and Bob alternately color an uncolored
vertex from a fixed palette while maintaining a proper coloring; Alice moves
first.  Alice wins if every vertex is eventually colored, and Bob wins if an
uncolored vertex has no available color.  The game chromatic number
$\gamechi(G)$ is the least palette size for which Alice has a winning
strategy.  Bodlaender's 1991 paper formulated graph-coloring games and
studied their algorithmic questions~\cite{Bodlaender1991}.

Mou, Sun, and Zhang use the generalized-Mycielski notation adopted here and
proved
\[
  4\leq \gamechi\bigl(M_k(C_3)\bigr)\leq 5
  \qquad (k\geq 2)
\]
\cite[Theorem~5.1(1)]{MouSunZhang2026}.  Their result leaves open which
endpoint occurs for the fixed two-layer instance.  For comparison, the
classical one-layer Mycielski graph of $C_3$ has game chromatic number four
\cite[p.~262]{AlagammaiVijayalakshmi2019}.  Our contribution is exactly the
two-layer endpoint; no claim of priority beyond the literature comparison
just stated is made.

Write $[3]=\{1,2,3\}$ and
\[
 V=\{v_i^m:i\in[3],\ m\in\{0,1,2\}\}\cup\{\omega\}.
\]
The graph studied below has edge set
\begin{align*}
E={}&\{v_i^0v_j^0:1\leq i<j\leq3\}\\
 &\cup\{v_i^m v_j^{m+1}:m\in\{0,1\},\ i,j\in[3],\ i\neq j\}\\
 &\cup\{v_i^2\omega:i\in[3]\}.
\end{align*}
Thus it has ten vertices and eighteen edges.  Only the base layer has
within-layer edges; consecutive layers are joined at unequal indices; and the
root is adjacent exactly to the three vertices in the top layer.

\begin{theorem}\label{thm:main}
Alice wins the four-color game on the graph $M_2(C_3)$ defined above.
Consequently,
\[
  \gamechi\bigl(M_2(C_3)\bigr)=4.
\]
\end{theorem}

\section{Lower bound}

The ordinary chromatic number is already at least four.  Indeed, in any
proper three-coloring the triangle $v_1^0v_2^0v_3^0$ uses all three colors.
For a fixed $j$, the vertex $v_j^1$ is adjacent to the two base vertices
$v_i^0$ with $i\neq j$, so it is forced to have the color of $v_j^0$.
Similarly, $v_j^2$ is forced to have the color of $v_j^1$.  The three
neighbors $v_1^2,v_2^2,v_3^2$ of $\omega$ therefore use all three colors,
leaving no color for $\omega$.  Hence
\[
  4\leq\chi\bigl(M_2(C_3)\bigr)
  \leq\gamechi\bigl(M_2(C_3)\bigr).
\]

\section{The exact strategy certificate}

We encode a partial coloring as a ten-tuple in the vertex order
\[
(v_1^0,v_2^0,v_3^0,v_1^1,v_2^1,v_3^1,
  v_1^2,v_2^2,v_3^2,\omega),
\]
using $0$ for uncolored and $1,2,3,4$ for colors.  The group
$S_3\times S_4$ acts by applying one index permutation simultaneously in all
three layers and by permuting the four color names.  Both actions preserve
the graph, legality, terminal status, and the player to move.  Each orbit is
represented by its lexicographically least tuple.

The file \texttt{strategy.json} contains the strategy directed acyclic graph.
At every nonterminal Alice node it records one legal vertex--color move.  At
every nonterminal Bob node it records the complete ordered list of all raw
legal vertex--color moves, before deduplication of symmetry-equivalent
successors.  The target of an edge is the canonical representative obtained
after making its move.  Terminal nodes are labeled either \texttt{all} or
\texttt{dead} according as the coloring is complete or an uncolored vertex
has no available color.

\begin{lemma}\label{lem:certificate}
Suppose a finite strategy graph has the preceding semantics, every Alice node
has a recorded legal move, every Bob node contains every legal reply, every
edge increases the number of colored vertices by one, and every reachable
leaf is labeled \texttt{all}.  Then the recorded strategy wins for Alice from
the initial state.
\end{lemma}

\begin{proof}
The number of uncolored vertices strictly decreases along every edge, so the
graph is well founded.  At a leaf Alice has already won.  Induct backward on
the number of uncolored vertices.  At an Alice node the recorded move reaches
a winning child; at a Bob node every move reaches a winning child because all
legal replies are present.  Symmetry does not change any of these properties,
so a move prescribed at a canonical representative transports to every state
in its orbit.  The initial empty state is therefore winning for Alice.
\end{proof}

The separately implemented program \path{verify_strategy.py} checks the
hypotheses of Lemma~\ref{lem:certificate} without importing the minimax solver.
It reconstructs the eighteen graph edges and explicitly enumerates all
$6\cdot24=144$ index--color actions to canonicalize each state.  It then
checks properness, move legality, terminal labels, turn parity, successor
states, the unique strategy move at each nonterminal Alice node, and exact
ordered equality between the listed and independently generated Bob moves.
The verified counts are
\[
\begin{array}{c|r}
\text{reachable canonical states} & 1020\\
\text{directed edges replayed} & 2467\\
\text{raw legal Bob replies covered} & 1809\\
\text{Alice strategy moves checked} & 658\\
\text{complete terminal states} & 36\\
\text{dead terminal states} & 0.
\end{array}
\]
Thus every hypothesis of Lemma~\ref{lem:certificate} holds, proving the
four-color upper bound in Theorem~\ref{thm:main}.  Together with the lower
bound from the preceding section, this proves the asserted equality.

The certificate architecture---a symmetry-reduced exhaustive object checked
by a reconstruction-oriented implementation---is also used in other exact
finite graph computations; for a recent example with independent orbit and
color-certificate reconstruction, see Zeng et
al.~\cite{ZengLiuRatnaveluOng2026}.  That paper is methodological context only;
the present proof depends on the certificate and verifier described above.

\section{Reproducibility and limitations}

From the project root, the decisive replay is
\begin{verbatim}
python3 evidence/verify_strategy.py evidence/strategy.json \
  --report evidence/verification.json
\end{verbatim}
The certificate can independently be regenerated by
\begin{verbatim}
python3 evidence/minimax_solver.py --output evidence/strategy.json
\end{verbatim}
Regeneration is not part of the proof's trust boundary: acceptance of the
upper bound rests on replay of the complete certificate.  All state and graph
operations are finite integer or symbolic operations; there are no random
choices, floating-point tolerances, or statistical inferences.

This theorem concerns only the fixed graph $M_2(C_3)$.  It does not determine
$\gamechi(M_k(C_3))$ for $k>2$, nor does it replace the general bounds in
\cite{MouSunZhang2026}.

\paragraph{Computational and AI assistance.}
Python programs performed the exact minimax search, produced the strategy
certificate, and replayed the certificate.  OpenAI Codex assisted with the
literature search, search-space and verifier design, adversarial audit
organization, and initial manuscript preparation.  Language-model output is
not mathematical evidence; the claimed upper bound is supported by the
explicit finite reduction and the independently implemented exact verifier.
No proof assistant was used.

\bibliographystyle{plain}
\bibliography{references}

\end{document}
